%%%%%%%%%%%%%%%%%%%%%%%%%%%%%%%%%%%%%%%%%%%%%%%%%%%%%%%%%%%%%%%%%%%%%%%%%%%%%
%                                                                           %
%                  COMMUTATORS AND ANALYTIC DEPENDENCE OF                   %
%                    FOURIER-BESSEL SERIES ON (0,\infty)                    %
%                                                                           %
%       COMIENZO ESCRITURA: 30 DE JUNIO DE 1997, POR JUAN LUIS VARONA.      %
%   ULTIMA MODIFICACION: 10 DE SEPTIEMBRE DE 1997, POR MARIO PEREZ RIERA.   %
%                                                                           %
%                  VERSION EN AMS-TEX 2.1, CON AMSPPT.STY                   %
%                                                                           %
%%%%%%%%%%%%%%%%%%%%%%%%%%%%%%%%%%%%%%%%%%%%%%%%%%%%%%%%%%%%%%%%%%%%%%%%%%%%%

\input amstex          % LINE NOT NECESSARY IF TYPESET WITH AMS-TEX 2.1
\documentstyle{amsppt} % NOT NEC. IF TYPESET WITH AMS-TEX PLUS AMSPPT.STY
\magnification=\magstep1
%\voffset=-0.8cm
%\nologo
\refstyle{C}

\define\N{{\Bbb N}}    % THE SET N OF NATURAL NUMBERS
\define\R{{\Bbb R}}    % THE SET R OF REAL NUMBERS

\define\BMO{\operatorname{BMO}} % THE SPACE BMO

\hyphenation{Pu-bli-ca-cions}

%%%%%%%%%%%%%%%%%%%%%%%%%%%%%%%%%%%%%%%%%%%%%%%%%%%%%%%%%%%%%%%%%%%%%%%%%%

\topmatter

\title Commutators and analytic dependence of Fourier-Bessel series on
$(0,\infty)$\endtitle

\rightheadtext{Commutators and analytic dependence of Fourier-Bessel series}

\author Jos\'e J.~Guadalupe, Mario P\'erez and	 Juan L.~Varona\endauthor

%\leftheadtext{J.~J.~Guadalupe, M.~P\'erez and J.~L.~Varona}

\address Jos\'e J.~Guadalupe. Dpto.~de Matem\'aticas y Computaci\'on.
Universidad de La Rioja. 26004 Logro\~no. Spain.\endaddress

\address Mario P\'erez. Dpto.~de Matem\'aticas. Universidad de Zaragoza.
50009 Zaragoza. Spain.\endaddress

\email mperez\@posta.unizar.es\endemail

\address Juan L.~Varona. Dpto.~de Matem\'aticas y Computaci\'on.
Universidad de La Rioja. 26004 Logro\~no. Spain.\endaddress

\email jvarona\@dmc.unirioja.es\endemail

\thanks Research supported by DGICYT and by UR.\endthanks

\thanks THIS PAPER HAS BEEN PUBLISHED IN:
{\it Canad.\ Math.\ Bull.}\ {\bf 42} (1999), {198--208}.
\endthanks

\keywords Fourier-Bessel series, commutators, $\BMO$, $A_p$
weights\endkeywords

\subjclass Primary 42C10\endsubjclass

\abstract In this paper we study the boundedness of the commutators
$[b, S_n]$ where $b$ is a $\BMO$ function and $S_n$ denotes the
$n$-th partial sum of the Fourier-Bessel series on $(0,\infty)$. Perturbing the
measure by $\exp(2b)$ we obtain that certain operators related to $S_n$ depend
analytically on the functional parameter $b$.
\endabstract

\endtopmatter

\document

%%%%%%%%%%%%%%%%%%%%%%%%%%%%%%%%%%%%%%%%%%%%%%%%%%%%%%%%%%%%%%%%%%%%%%%%%%%

\head 0. Introduction.\endhead

Let $J_\alpha$ be the Bessel function of order $\alpha > -1$. The formula
$$
   \int_0^\infty J_{\alpha + 2 n + 1}(x) J_{\alpha + 2 m + 1}(x) \,
   {dx \over x} =
   \cases 0,                             &{\text{ if }} n \not= m\cr
          2^{-1} (\alpha + 2n + 1)^{-1}, &{\text{ if }} n = m
    \endcases
$$
(see~\cite{14, XIII.13.41 (7), p.~404} and~\cite{14, XIII.13.42 (1), p.~405})
provides an orthonormal system $(j_n^\alpha)_{n\geq 0}$ in $L^2((0,\infty),
x^\alpha \, dx)$ [$L^2(x^\alpha)$, from now on], given by
$$
   j_n^\alpha(x) = \sqrt{\alpha + 2n + 1} J_{\alpha + 2n + 1}(\sqrt x)
   \, x^{-\alpha/2 - 1/2}.
$$
In this paper we consider the Fourier expansion associated with this
orthonormal system, which is usually referred to as the Fourier-Bessel series
on $(0,\infty)$. For any suitable function $f$ and any $n \geq 0$, the $n$-th
partial sum of this expansion is given by
$$
   S_n f = \sum_{k=0}^n c_k(f) j_k^\alpha,
   \qquad c_k(f) = \int_0^\infty f(t) j_k^\alpha(t) t^\alpha \, dt.
$$
We also consider the commutator of the Fourier-Bessel series on $(0,\infty)$
and the multiplication operator associated to a $\BMO$ function; this is
defined, for any given $b \in \BMO$ and $n \geq 0$, as
$$
   [b, S_n] f = b S_n(f) - S_n(bf).
$$
In the case $\alpha \geq -1/2$, one of the authors proved in~\cite{13} that the
Fourier-Bessel series is bounded in $L^p(x^\alpha)$, i.e., there exists some
constant $C > 0$ (depending on $\alpha$ and $p$) such that for every $n \geq
0$ and every $f \in L^p(x^\alpha)$,
$$
   \| S_n f \|_{L^p(x^\alpha)} \leq C \| f \|_{L^p(x^\alpha)},
$$
if and only if $\max\{4/3, 4(\alpha+1)/(2\alpha+3)\} < p < \min\{4,
4(\alpha+1)/(2\alpha+1)\}$. In Theorem~1 we will extend this result to the case
$\alpha > -1$ and prove the corresponding inequality for the commutator $[b,
S_n]$, $b \in \BMO$.

Regarding the commutator $[b, S_n]$, results of this type are of independent
interest and have been widely studied for many classical operators;
see~\cite{2, 10, 11, 12, 4}, for instance.

In our case, the commutator $[b, S_n]$ is closely related to the problem of
perturbating the orthonormal system. Given an orthonormal system
$(\varphi_n)_{n \geq 0}$ in some $L^2(\nu)$ space and a suitable function $b$
(in some sense close to $0$), the classical Gram-Schmidt procedure can be
applied to $(\varphi_n)_{n \geq 0}$ so as to obtain a new orthonormal system
in $L^2(e^{2b} d\nu)$, which we will refer to as a perturbated system. In this
natural way a mapping can be defined that associates a perturbated system (and
a perturbated orthogonal expansion) to each (small) function $b$. For different
compact perturbations of orthogonal polynomial systems and further references,
see~\cite{7,~9,~1}.

Let us take the system $(j_n^\alpha)_{n \geq 0}$ in $L^2(x^\alpha)$ as our
starting point. Let ${\bold S}_n(b)$ stand for the $n$-th partial sum operator
of the Fourier series associated to the perturbed measure $e^{2b} x^\alpha \,
dx$ in the aforementioned way. Once the boundedness properties of $S_n =
{\bold S}_n(0)$ have been established, it is interesting to study the mapping
$b
\mapsto {\bold S}_n(b)$. This is not, however, a convenient setting, since each
perturbed series
${\bold S}_n(b)$ acts on a different space $L^2(e^{2b} x^\alpha)$. Instead, we
can consider the operators
$$
   V_n(b) = e^b {\bold S}_n(b) e^{-b}.
$$
Now, each $V_n(b)$ acts on $L^2(x^\alpha)$ and its norm coincides with the
operator norm of ${\bold S}_n(b)$ acting on $L^2(e^{2b} x^\alpha)$. The problem
is further simplified if we take the operators
$$
   T_n(b) = e^b {\bold S}_n(0) e^{-b},
$$
i.e., $T_n(b)f = e^b S_n(e^{-b}f)$. Indeed, it has been proved in~\cite{3}
that the family $(V_n(b))_{n\geq 0}$ depends analytically on $b$ belonging to a
neighbourhood of $0$ in the complexification of $\BMO$ whenever the family
$(T_n(b))_{n\geq 0}$ does too.

We will prove in Theorem~2 that the family of operators $(T_n(b))_{n \geq 0}$
acting on $L^2(x^\alpha)$ is uniformly bounded for $b$ belonging to some
neighbourhood of $0$ in the complexification of $\BMO$. As a consequence
(see~\cite{3, Propositions 2.1 and 2.3}), the operator-valued mappings
$(T_n)_{n \geq 0}$ are uniformly analytic in a neighbourhood of $0$ in the
complexification of $\BMO$ and so are $(V_n)_{n \geq 0}$.

Now, the connection between $[b, S_n]$ and the perturbated Fourier series
comes via the G\^ateaux differential of $T_n$ at $0$ in the direction $b$:
$$
   {d \over dz} T_n(zb) \big|_{z=0} = [b,S_n].
$$
In this way, the uniform analyticity of $T_n$ in a neighbourhood of $0$ gives
the $L^2$-boundedness of $[b, S_n]$.

%%%%%%%%%%%%%%%%%%%%%%%%%%%%%%%%%%%%%%%%%%%%%%%%%%%%%%%%%%%%%%%%%%%%%%%%%%

\head 1. Main results.\endhead

If $b$ is a locally Lebesgue-integrable function on $(0, \infty)$, the mean of
$b$ on an interval $I \subseteq (0,\infty)$ is
$$
   b_I = {1 \over |I|} \int_I b(x) \, dx.
$$
The function $b$ is said to have bounded mean oscillation on $(0,\infty)$ if
$$
   \| b \|_{\BMO} = \sup_I {1 \over |I|} \int_I |b(x) - b_I| \, dx
$$
is finite, where the supremum is taken over all the intervals $I \subseteq
(0,\infty)$. The space $\BMO$ of real-valued functions (modulo constants)
having bounded mean oscillation on $(0,\infty)$ is a real Banach space with $\|
\cdot \|_{\BMO}$ as its norm.

\proclaim{Theorem 1} Let $1 < p < \infty$, $-1 < \alpha$ such that
$$
   \cases
   4/3 < p < 4, &{\text{ if }} \ -1 < \alpha < 0; \cr
   \displaystyle
   {4 (\alpha + 1) \over 2 \alpha + 3} < p < {4(\alpha+1) \over 2 \alpha + 1},
                &{\text{ if }} \ 0 \leq \alpha.
   \endcases
$$
\roster
\item"(a)" There exists some constant $C > 0$ such that, for every $f \in
L^p(x^\alpha)$ and $n \geq 0$,
$$
   \| S_n f \|_{L^p(x^\alpha)} \leq C \| f \|_{L^p(x^\alpha)}.
$$
\item"(b)" If $b \in \BMO$, then there exists some constant $C > 0$ such that,
for every $f \in L^p(x^\alpha)$ and $n \geq 0$,
$$
   \| [S_n, b] f \|_{L^p(x^\alpha)} \leq C \| f \|_{L^p(x^\alpha)}.
$$
\endroster
\endproclaim

Throughout this paper, we will denote by $C$ a positive constant which is
independent of $n$ and $f$, but may be different in each occurrence, even
within the same formula.

\proclaim{Theorem 2}  Let $1 < p < \infty$, $-1 < \alpha$ such that
$$
   \cases
   4/3 < p < 4, &{\text{ if }} \ -1 < \alpha < 0; \cr
   \displaystyle
   {4 (\alpha + 1) \over 2 \alpha + 3} < p < {4(\alpha+1) \over 2 \alpha + 1},
                &{\text{ if }} \ 0 \leq \alpha.
   \endcases
$$
Then there exist some $C, \delta > 0$ such that, for all $b \in \BMO$
with $\| b \|_{\BMO} < \delta$,
$$
   \sup_n \| T_n(b) \|_{L^p(x^\alpha) \to L^p(x^\alpha)} \leq C.
$$
\endproclaim

The next corollary is just a consequence of Theorem~2
and~\cite{3, Prop.~2.3}.

\proclaim{Corollary} The sequences of operators $(T_n(b))_{n \geq 0}$ and
$(V_n(b))_{n \geq 0}$, acting on the space $L^2(x^\alpha)$, are uniformly
analytic in a neighbourhood of $0$ in the complexification of $\BMO$.
\endproclaim

Some notation and previous results will be necessary. For $1 < p < \infty$, we
write $p' = p/(p-1)$, i.e., $1/p + 1/p' = 1$. A weight is a nonnegative
Lebesgue-measurable function on $(0,\infty)$. The class $A_p(0,\infty)$
[$A_p$, for short] consists of those pairs of weights $(u,v)$ such that, for
every subinterval $I \subseteq (0, \infty)$,
$$
   {1 \over |I|} \int_I u
   \left( {1 \over |I|} \int_I v^{-p'/p} \right)^{p/p'} \leq C,
$$
where $C$ is a positive constant independent of $I$, and $|I|$ denotes the
length of $I$. The $A_p$ constant of $(u,v)$ is the least constant $C$
verifying this inequality and will be denoted by $A_p(u,v)$. A single weight
$w$ is said to belong to $A_p$ if $(w,w) \in A_p$; in this case we denote the
constant by $A_p(w)$. We refer the reader to~\cite{6} for further details on
$A_p$ classes.

The Hilbert transform on $(0,\infty)$ will be denoted by $H$. Fix $1 < p <
\infty$; then $H$ is a bounded linear operator on $L^p(w)$, for any weight $w
\in A_p$. The norm of $H \colon L^p(w) \longrightarrow L^p(w)$ and the $A_p$
constant of $w$ depend only one on another, in the sense that given some
constant $C$ which verifies the $A_p$ condition for $w$, another constant
$C_1$ depending only on $C$ can be chosen so that $\| H\| \leq C_1$, and
viceversa. Therefore, for a sequence $(w_n)_{n \in \N}$ uniformly in $A_p$,
i.e., with some constant $C$ verifying the $A_p$ condition for every $w_n$, the
Hilbert transform is uniformly bounded on $L^p(w_n)$, $n \in \N$. We refer the
reader again to~\cite{6} for further details.

Also, if $(u, v)$ is a pair of weights such that $C_1 u \leq w \leq C_2 v$ for
some $w \in A_p$, we deduce that $H$ is a bounded operator from $L^p(v)$ into
$L^p(u)$. The existence of such a weight $w$ is equivalent to $(u^\delta,
v^\delta) \in A_p$ for some $\delta > 1$ (see~\cite{8}). For short, this is
written as $(u,v) \in A_p^\delta$.

Analogous results hold also with the commutator $[b, H]$, for any $b \in \BMO$
(see~\cite{2}, for instance). Namely, given $b \in \BMO$ and $w \in A_p$, $[b,
H]$ is a bounded operator on $L^p(w)$ with a norm that depends only on the
$\BMO$-norm of $b$ and the $A_p$ constant of $w$, in the sense above.

%%%%%%%%%%%%%%%%%%%%%%%%%%%%%%%%%%%%%%%%%%%%%%%%%%%%%%%%%%%%%%%%%%%%%%%%%%

\head 2. Proofs.\endhead

Let us start with some auxiliary results:

\proclaim{Lemma 1} Let $u$, $v$, $w$ be weights on $(0, +\infty)$,
$\lambda > 0$.
\roster
\item"(a)" $w(x) \in A_p$ if and only if $w(\lambda x) \in A_p$;
both weights have the same $A_p$ constant.
\item"(b)" $w \in A_p$ if and only if $\lambda w \in A_p$; both weights
have also the same $A_p$ constant.
\item"(c)" If $u, v \in A_p$, then $u + v \in A_p$ and
$A_p(u + v) \leq A_p(u) + A_p(v)$.
\item"(d)" If $u, v \in A_p$ and $1/w = 1/u + 1/v$, then $w \in A_p$
and $A_p(w) \leq C [A_p(u) + A_p(v)]$.
\endroster
\endproclaim

\demo{Proof} Parts (a) and (b) are trivial. Part (c) follows easily from
the inequality
$$
   \left(
   {1 \over |I|} \int_I (u + v)^{-{p' / p}} \right)^{p / p'}
   \leq
   \min\left\{
   \left({1 \over |I|} \int_I u^{-{p' / p}} \right)^{p / p'},
   \left({1 \over |I|} \int_I v^{-{p' / p}} \right)^{p / p'}
   \right\}.
$$
Part (d) is a consequence of (c) and the fact that $u \in A_p \iff
u^{-p'/p} \in A_{p'}$, with $A_{p'}(u^{-p'/p}) = [A_p(u)]^{p'/p}$. \qed
\enddemo

The proof of the next lemma is not difficult, but cumbersome, so we omit it.
For the weight in (c), observe that $x^r |x^{1/2} - 1|^s \sim x^r$ near
$0$, $x^r |x^{1/2} - 1|^s \sim |x-1|^s$ near $1$ and $x^r |x^{1/2} - 1|^s \sim
x^{r+s/2}$ near $\infty$, whence the three conditions follow.

\proclaim{Lemma 2} Let $r, s \in \R$.
\roster
\item"(a)" $x^r \in A_p \iff -1 < r < p-1$.
\item"(b)" Set $\Phi(x) = x^r$ if $x \in (0,1)$ and $\Phi(x) = x^s$ if $x
\in (1,\infty)$. Then, $\Phi \in A_p$ if and only if $-1 < r < p-1$ and $-1 <
s < p-1$.
\item"(c)" $x^r |x^{1/2} - 1|^s \in A_p \iff -1 < r < p-1, \ -1 < s < p-1$ and
$-1 < r + s/2 < p-1$.
\endroster
\endproclaim

\proclaim{Lemma 3} Let $n \in \N$, $\alpha > -1$. Then
$$
   \eqalign{
   \sum_{k=0}^n 2 (\alpha + 2k + 1) &J_{\alpha + 2k + 1}(x) J_{\alpha + 2k +
   1}(t)
\cr
   = {x t \over x^2 - t^2} \, [
   &x J_{\alpha+1}(x) J_\alpha(t) - t J_\alpha(x) J_{\alpha+1}(t) \cr
   &+ x J'_{\alpha + 2n + 2}(x) J_{\alpha + 2n + 2}(t)
   - t J_{\alpha + 2n + 2}(x) J'_{\alpha + 2n + 2}(t)
   ].
   }
$$
\endproclaim

\demo{Proof} Using the equality $J_{\nu-1}(z) + J_{\nu+1}(z) = {2 \nu
\over z} J_\nu(z)$ (see~\cite{14, III.3.2, p.~45}) to express $J_{\mu-1}$
and $J_{\mu+2}$ in terms of $J_\mu$ and $J_{\mu+1}$ proves the formula
$$
  {x t \over x^2 - t^2} \left[
   x J_\mu(x) J_{\mu-1}(t) - t J_{\mu-1}(x) J_\mu(t)
   - x J_{\mu + 2}(x) J_{\mu + 1}(t)
   + t J_{\mu + 1}(x) J_{\mu + 2}(t)
   \right]
$$
$$
   = 2 \mu J_\mu(x) J_\mu(t).
$$
This gives now
$$
   \eqalign{
   \sum_{k=0}^n
   2 (\alpha + 2k + 1) &J_{\alpha + 2k + 1}(x) J_{\alpha + 2k + 1}(t)
   \cr
   = {x t \over x^2 - t^2} \, [
   &x J_{\alpha+1}(x) J_\alpha(t) - t J_\alpha(x) J_{\alpha+1}(t)
   \cr
   &- x J_{\alpha + 2n + 3}(x) J_{\alpha + 2n + 2}(t)
   + t J_{\alpha + 2n + 2}(x) J_{\alpha + 2n + 3}(t)
   ].
   }
$$
Finally, use the formula $z J_{\nu+1}(z) = \nu J_\nu(z) - z J'_\nu(z)$
(see~\cite{14, III.3.2, p.~45}) to take out $J_{\alpha+2n+3}$.
\qed
\enddemo

\demo{Proof of Theorem 1} From the definition,
$$
   S_nf(x) = x^{-{\alpha \over 2} - {1 \over 2}} \int_0^\infty
   \left[ \sum_{k=0}^n (\alpha + 2k + 1) J_{\alpha + 2k + 1}(x^{1 \over 2})
   J_{\alpha + 2k + 1}(t^{1 \over 2}) \right]
   t^{{\alpha \over 2} - {1 \over 2}} f(t) \, dt
$$
so that Lemma~3 leads to
$$
   S_n f = W_1 f - W_2 f + W_{3,n} f - W_{4,n} f,
$$
where
$$
   \eqalign{
   W_1 f (x) &= 2^{-1} x^{-{\alpha / 2} + {1 / 2}} 
   J_{\alpha+1}(x^{1 / 2})
   \, H\big(t^{\alpha / 2} J_\alpha(t^{1 / 2}) f(t) \big) (x),
   \cr
   W_2 f (x) &= 2^{-1} x^{-{\alpha / 2}} 
   J_\alpha(x^{1 / 2}) \,
   H\big(t^{{\alpha / 2} + {1 / 2}} J_{\alpha+1}(t^{1 / 2})
   f(t) \big) (x),
   \cr
   W_{3,n} f (x) &= 2^{-1} x^{-{\alpha / 2} + {1 / 2}}
   J_\nu'(x^{1 / 2})
   \, H\big(t^{\alpha / 2} J_\nu(t^{1 / 2}) f(t)\big) (x),
   \cr
   W_{4,n} f (x) &= 2^{-1} x^{-{\alpha / 2}} 
   J_\nu(x^{1 / 2}) \,
   H\big(t^{{\alpha / 2} + {1 / 2}} J_\nu'(t^{1 / 2}) f(t)\big)(x)
   }
$$
and $\nu = \alpha + 2n + 2$. Thus, we will show that the operators $W_1$,
$W_2$ are bounded and the operators $W_{3,n}$, $W_{4,n}$ are uniformly bounded
for $n \geq 0$. The proof for the commutator $[b, S_n]$ is the same: just
put $[b, H]$ instead of $H$.

(I) {\it Boundedness of the operator $W_1$.} From the definition, it follows
that
$$
   \| W_1 f \|_{L^p(x^\alpha)} \leq C \| f \|_{L^p(x^\alpha)}
$$
if and only if
$$
   \| H g \|_{L^p(x^{\alpha - {\alpha p / 2} + {p / 2}}
   |J_{\alpha+1}(x^{1/2})|^p)} \leq C
   \| g \|_{L^p(x^{\alpha - {\alpha p / 2}} |J_\alpha(x^{1/2})|^{-p})}.
$$
Proving that there is a weight $\Phi \in A_p$ with
$$
   C x^{\alpha - {\alpha p / 2} + {p / 2}} |J_{\alpha+1}(x^{1/2})|^p
   \leq \Phi(x) \leq
   C x^{\alpha - {\alpha p / 2}} |J_\alpha(x^{1/2})|^{-p}
   \tag 1
$$
will be enough. According to the bounds
$$
   \eqalign{
   |J_\alpha(x)| &\leq C_\alpha x^\alpha, \quad x \in (0,1),
   \cr
   |J_\alpha(x)| &\leq C_\alpha x^{-1/2}, \quad x \in (1,\infty)
   \cr
   }
$$
(see, e.g.,~\cite{14, III.3.1 (8), p.~40} and~\cite{14, VII.7.21 (1),
p.~199}), we have
$$
   \eqalign{
   x^{\alpha - {\alpha p / 2} + {p / 2}} |J_{\alpha+1}(x^{1/2})|^p
   &\leq
   \cases
   C x^{\alpha + p}, &{\text{ if }} x \in (0,1), \cr
   C x^{\alpha - {\alpha p / 2} + {p / 4}},
   &{\text{ if }} x \in (1,\infty),
   \endcases   
   \cr
   x^{\alpha - {\alpha p / 2}} |J_\alpha(x^{1/2})|^{-p}
   &\geq
   \cases
   C x^{\alpha - \alpha p}, &{\text{ if }} x \in (0,1), \cr
   C x^{\alpha - {\alpha p / 2} + {p / 4}},
   &{\text{ if }} x \in (1,\infty).
   \endcases
   }
$$
Let us try
$$
   \Phi(x) = \cases
   x^r, &{\text{ if }} x \in (0,1), \cr
   x^{\alpha - {\alpha p / 2} + {p / 4}},
   &{\text{ if }} x \in (1,\infty).
   \endcases
$$
By (b) in Lemma~2, conditions (1) and $\Phi \in A_p$ will hold if
$$
   \cases
   \alpha - \alpha p \leq r \leq \alpha + p, \cr
   -1 < r < p - 1, \cr
   -1 < \alpha - {\alpha p / 2} + {p / 4} < p - 1.
   \endcases
$$
The third line is equivalent to
$$
   {2 \alpha - 1 \over 4} \, p < \alpha + 1,
   \qquad
   {\alpha + 1} < {2 \alpha + 3 \over 4} \, p,
$$
and these hold, by the hypothesis. For the inequalities involving $r$ it
suffices
$$
   \max\{-1, \alpha - \alpha p\} < \min\{p - 1, \alpha + p\}.
$$
It is easy to check that this also holds, whenever $\alpha > -1$ and $p > 1$.

(II) {\it Boundedness of the operator $W_2$.} The proof is entirely similar: we
have
$$
   \| W_2 f \|_{L^p(x^\alpha)} \leq C \| f \|_{L^p(x^\alpha)}
$$
if and only if
$$
   \| H g \|_{L^p(x^{\alpha - {\alpha p / 2}} |J_\alpha(x^{1/2})|^p)}
   \leq C
   \| g \|_{L^p(x^{\alpha - {\alpha p / 2} - {p / 2}}
   |J_{\alpha+1}(x^{1/2})|^{-p})}
$$
so that we can prove that there is a weight $\Psi \in A_p$ with
$$
   C x^{\alpha - {\alpha p / 2}} |J_\alpha(x^{1/2})|^p
   \leq \Psi(x) \leq
   C x^{\alpha - {\alpha p / 2} - {p / 2}}
   |J_{\alpha+1}(x^{1/2})|^{-p}.
   \tag 2
$$
Now we have
$$
   \eqalign{
   x^{\alpha - {\alpha p / 2}} |J_\alpha(x^{1/2})|^p
   &\leq
   \cases
   C x^\alpha, &{\text{ if }} x \in (0,1), \cr
   C x^{\alpha - {\alpha p / 2} - {p / 4}},
   &{\text{ if }} x \in (1,\infty),
   \endcases   
   \cr
   x^{\alpha - {\alpha p / 2} - {p / 2}} |J_{\alpha+1}(x^{1/2})|^{-p}
   &\geq
   \cases
   C x^{\alpha - \alpha p - p}, &{\text{ if }} x \in (0,1), \cr
   C x^{\alpha - {\alpha p / 2} - {p / 4}},
   &{\text{ if }} x \in (1,\infty).
   \endcases
   }
$$
Setting
$$
   \Psi(x) = \cases
   x^r, &{\text{ if }} x \in (0,1), \cr
   x^{\alpha - {\alpha p / 2} - {p / 4}},
   &{\text{ if }} x \in (1,\infty),
   \endcases
$$
conditions (2) and $\Psi \in A_p$ will hold if
$$
   \cases
   \alpha - \alpha p - p \leq r \leq \alpha, \cr
   -1 < r < p - 1, \cr
   -1 < \alpha - {\alpha p / 2} - {p / 4} < p - 1.
   \endcases
$$
The third line is equivalent to
$$
   {2 \alpha + 1 \over 4} \, p < \alpha + 1,
   \qquad
   {\alpha + 1} < {2 \alpha + 5 \over 4} \, p,
$$
and these hold, by the hypothesis. For the inequalities involving $r$ we
only need
$$
   \max\{-1, \alpha - \alpha p - p\} < \min\{p - 1, \alpha\}.
$$
It is easy to check that this also holds, whenever $\alpha > -1$ and $p > 1$.

(III) {\it Uniform boundedness of the operators $W_{3,n}$.} Here,
$$
   \| W_{3,n} f \|_{L^p(x^\alpha)} \leq C \| f \|_{L^p(x^\alpha)}
$$
if and only if
$$
   \| H g \|_{L^p(x^{\alpha - {\alpha p / 2} + {p / 2}}
   |J_\nu'(x^{1/2})|^p)} \leq C
   \| g \|_{L^p(x^{\alpha - {\alpha p / 2}} |J_\nu(x^{1/2})|^{-p})}.
$$
We make now use of the bounds
$$
   \eqalign{
   |J_\nu(x)| &\leq C x^{-1/4} \left[ |x - \nu| + \nu^{1/3} \right]^{-1/4},
   \quad \nu = \alpha + 2n + 2, \ x \in (0,\infty),
   \cr
   |J_\nu'(x)| &\leq C x^{-3/4} \left[ |x - \nu| + \nu^{1/3} \right]^{1/4},
   \quad \nu = \alpha + 2n + 2, \ x \in (0,\infty),
   }
$$
with some universal constant $C$. They follow from those of~\cite{5}, for
instance. Therefore,
$$
   \eqalign{
   x^{\alpha - {\alpha p / 2} + {p / 2}} |J_\nu'(x^{1/2})|^p
   &\leq C x^{\alpha - {\alpha p / 2} + {p / 8}}
   \left[|x^{1/2} - \nu| + \nu^{1/3} \right]^{p/4},
   \cr
   x^{\alpha - {\alpha p / 2}} |J_\nu(x^{1/2})|^{-p}
   &\geq C x^{\alpha - {\alpha p / 2} + {p / 8}}
   \left[|x^{1/2} - \nu| + \nu^{1/3} \right]^{p/4}.
   }
$$
It will suffice to prove that $\varphi_\nu \in A_p$ uniformly in $n$, with
$$
   \varphi_\nu(x) = x^{\alpha - {\alpha p / 2} + {p / 8}}
   \left[|x^{1/2} - \nu| + \nu^{1/3} \right]^{p/4}.
   \tag 3
$$
{}From Lemma~1, we have
$$
   \eqalign{
   \varphi_\nu(x) \in A_p {\text{ unif.}}
   &\iff \varphi_\nu(\nu^2 x) \in A_p {\text{ unif.}}
   \cr
   &\iff x^{\alpha - {\alpha p / 2} + {p / 8}}
   \left[ |x^{1/2} - 1| + \nu^{-2/3} \right]^{p/4} \in A_p {\text{ unif.}}
   \cr
   &\iff x^{\alpha - {\alpha p / 2} + {p / 8}} |x^{1/2} - 1|^{p/4}
   + \nu^{-p/6} x^{\alpha - {\alpha p / 2} + {p / 8}}
   \in A_p {\text{ unif.}},
   }
$$
where the last equivalence follows from
$$
   \left[ |x^{1/2} - 1| + \nu^{-2/3} \right]^{p/4}
   \sim |x^{1/2} - 1|^{p/4} + \nu^{-p/6},
$$
i.e., the ratio of both terms is bounded below and above by two positive
constants not depending on $n$ or $x$. Now, again by Lemma~1, proving that
$x^{\alpha - {\alpha p / 2} + {p / 8}} \in A_p$
and $x^{\alpha - {\alpha p / 2} + {p / 8}} |x^{1/2} - 1|^{p/4} \in A_p$
will suffice. According to Lemma~2,
$$
   \eqalign{
   \cases
   x^{\alpha - {\alpha p / 2} + {p / 8}} \in A_p \cr
   x^{\alpha - {\alpha p / 2} + {p / 8}} |x^{1/2} - 1|^{p/4} \in A_p
   \endcases
   &\iff \cases
   -1 < \alpha - {\alpha p / 2} + {p / 8} < p - 1 \cr
   -1 < {p / 4} < p - 1 \cr
   -1 < \alpha - {\alpha p / 2} + {p / 4} < p - 1
   \endcases
   \cr
   &\iff \cases
   -1 < \alpha - {\alpha p / 2} + {p / 8} \cr
   4/3 < p \cr
   \alpha - {\alpha p / 2} + {p / 4} < p - 1
   \endcases
   \cr
   &\iff \cases
   {2 \alpha - 1/2 \over 4} \, p < \alpha + 1 < {2 \alpha + 3 \over 4} \, p \cr
   4/3 < p
   \endcases
   }
$$
and these inequalities follow from the initial conditions.

(IV) {\it Uniform boundedness of the operators $W_{4,n}$.} Finally,
$$
   \| W_{4,n} f \|_{L^p(x^\alpha)} \leq C \| f \|_{L^p(x^\alpha)}
$$
if and only if
$$
   \| H g \|_{L^p(x^{\alpha - {\alpha p / 2}} |J_\nu(x^{1/2})|^p)}
   \leq C
   \| g \|_{L^p(x^{\alpha - {\alpha p / 2} - {p / 2}}
   |J_\nu'(x^{1/2})|^{-p})}.
$$
Also,
$$
   \eqalign{
   x^{\alpha - {\alpha p / 2}} |J_\nu(x^{1/2})|^p
   &\leq C x^{\alpha - {\alpha p / 2} - {p / 8}}
   \left[|x^{1/2} - \nu| + \nu^{1/3} \right]^{-p/4},
   \cr
   x^{\alpha - {\alpha p / 2} - {p / 2}} |J_\nu'(x^{1/2})|^{-p}
   &\geq C x^{\alpha - {\alpha p / 2} - {p / 8}}
   \left[|x^{1/2} - \nu| + \nu^{1/3} \right]^{-p/4}
   }
$$
so let us put
$$
   \psi_\nu(x) = x^{\alpha - {\alpha p / 2} - {p / 8}}
   \left[ |x^{1/2} - \nu| + \nu^{1/3} \right]^{-p/4}
   \tag 4
$$
and show that $\psi_\nu \in A_p$ uniformly in $n$. Indeed,
$$
   \eqalign{
   \psi_\nu(x) \in A_p {\text{ unif.}}
   &\iff \psi_\nu(\nu^2 x) \in A_p {\text{ unif.}}
   \cr
   &\iff x^{\alpha - {\alpha p / 2} - {p / 8}}
   \left[ |x^{1/2} - 1| + \nu^{-2/3} \right]^{-p/4} \in A_p {\text{ unif.}}
   }
$$
and
$$
   \eqalign{
   \Big(x^{\alpha - {\alpha p / 2} - {p / 8}}
   &\left[ |x^{1/2} - 1| + \nu^{-2/3} \right]^{-p/4}
   \Big)^{-1}
   \cr
   &\sim x^{-\alpha + {\alpha p / 2} + {p / 8}}
   \left[|x^{1/2} - 1|^{p/4} + \nu^{-p/6} \right]
   \cr
   &= \left[x^{\alpha - {\alpha p / 2} - {p / 8}}
   |x^{1/2} - 1|^{-p/4} \right]^{-1}
   + \left[ \nu^{p/6} x^{\alpha - {\alpha p / 2} - {p / 8}}
   \right]^{-1}
   }
$$
so that proving that $x^{\alpha - {\alpha p / 2} - {p / 8}}
|x^{1/2} - 1|^{-p/4} \in A_p$ and $x^{\alpha - {\alpha p / 2} - {p / 8}} \in
A_p$ will suffice. But
$$
   \eqalign{
   \cases
   x^{\alpha - {\alpha p / 2} - {p / 8}} \in A_p \cr
   x^{\alpha - {\alpha p / 2} - {p / 8}} |x^{1/2} - 1|^{-p/4} \in A_p
   \endcases
   &\iff \cases
   -1 < \alpha - {\alpha p / 2} - {p / 8} < p - 1 \cr
   -1 < - {p / 4} < p - 1 \cr
   -1 < \alpha - {\alpha p / 2} - {p / 4} < p - 1
   \endcases
   \cr
   &\iff \cases
   \alpha - {\alpha p / 2} - {p / 8} < p - 1 \cr
   p < 4 \cr
   -1 < \alpha - {\alpha p / 2} - {p / 4}
   \endcases
   \cr
   &\iff \cases
   {2 \alpha + 1 \over 4} \, p < \alpha + 1 < {2 \alpha + 9/2 \over 4} \, p
   \cr
   p < 4
   \endcases
   }
$$
and these inequalities hold by the hypothesis. The proof of Theorem~1 is now
complete.
\qed
\enddemo

\demo{Proof of Theorem 2} For each $n \geq 0$ and $b \in \BMO$, $T_n(b) \colon
L^p(x^\alpha) \longrightarrow L^p(x^\alpha)$ is bounded if and only if $S_n
\colon L^p(e^{pb} x^\alpha) \longrightarrow L^p(e^{pb} x^\alpha)$ is bounded,
and both operators have the same norm. Thus, we can follow the proof of
Theorem~1 and conclude that conditions (1), (2), (3) and (4), i.e.,
$$
   \eqalign{
   C x^{\alpha - {\alpha p / 2} + {p / 2}} |J_{\alpha+1}(x^{1/2})|^p
   \leq \Phi(x) &\leq
   C x^{\alpha - {\alpha p / 2}} |J_\alpha(x^{1/2})|^{-p},
   \cr
   C x^{\alpha - {\alpha p / 2}} |J_\alpha(x^{1/2})|^p
   \leq \Psi(x) &\leq
   C x^{\alpha - {\alpha p / 2} - {p / 2}}
   |J_{\alpha+1}(x^{1/2})|^{-p},
   \cr
   \varphi_\nu(x) &= x^{\alpha - {\alpha p / 2} + {p / 8}}
   \left[|x^{1/2} - \nu| + \nu^{1/3} \right]^{p/4},
   \cr
   \psi_\nu(x) &= x^{\alpha - {\alpha p / 2} - {p / 8}}
   \left[|x^{1/2} - \nu| + \nu^{1/3} \right]^{-p/4},
   }
$$
are still sufficient, if we require now $e^{pb} \Phi, e^{pb} \Psi, e^{pb}
\varphi_\nu, e^{pb} \psi_\nu \in A_p$ uniformly in $\nu$. The proof of
Theorem~1, together with next lemma, finish the proof of Theorem~2.
\qed
\enddemo

\proclaim{Lemma 4} Let $1 < p < \infty$. For each $\phi \in A_p$, there exists
some $\delta > 0$ such that $e^{p b} \phi \in A_p$ whenever $b \in \BMO$ with
$\| b \|_{\BMO} < \delta$. Moreover, $\delta$ and the $A_p$ constant of $e^{p
b} \phi$ depend only on the $A_p$ constant of $\phi$.
\endproclaim

\remark{Remark} Again, statements like ``$\delta$ depends only on the $A_p$
constant of $\phi$'' should be understood as: given a constant $C > 0$ which
verifies the $A_p$ condition for $\phi$, some $\delta$ can be chosen
depending only on $C$.
\endremark

\demo{Proof} If $\phi \in A_p$, there exists some $\varepsilon > 1$
such that $\phi^\varepsilon \in A_p$; moreover, $\varepsilon$ and the $A_p$
constant of $\phi^\varepsilon$ depend only on the $A_p$ constant of
$\phi$~\cite{6, Theorem~IV.2.7, p.~399}. Take now $1/\varepsilon +
1/\varepsilon' = 1$. There exists some $\delta > 0$ such that
$$
   \| b \|_{\BMO} < \delta \Longrightarrow e^{p \varepsilon' b} \in A_p;
$$
here, $\delta$ and the $A_p$ constant of $e^{p \varepsilon' b}$ depend only on
$\varepsilon'$~\cite{6, p.~409}. This, together with $\phi^\varepsilon \in
A_p$ and H\"older's inequality, imply $e^{p b} \phi \in A_p$ with an $A_p$
constant depending only on the $A_p$ constant of $\phi$.
\qed
\enddemo

%%%%%%%%%%%%%%%%%%%%%%%%%%%%%%%%%%%%%%%%%%%%%%%%%%%%%%%%%%%%%%%%%%%%%%%%%%
%%%%%%%%%%%%%%%%%%%%%%%%%%%%%%%%%%%%%%%%%%%%%%%%%%%%%%%%%%%%%%%%%%%%%%%%%%

\Refs

\widestnumber\key{14}

\ref
\key   1
\by    A. I. Aptekarev, A. Branquinho and F. Marcell\'an
\paper Toda-type differential equations for the recurrence coefficients of
       orthogonal polynomials and Freud transformation
\jour  J. Comput. Appl. Math.
\vol 78
\yr 1997
\pages 139--160
\endref

\ref
\key   2
\by    S. Bloom
\paper A commutator theorem and weighted $\BMO$
\jour  Trans. Amer. Math. Soc.
\vol   292
\yr    1985
\pages 103--122
\endref

\ref
\key   3
\by    R. R. Coifman and M. A. M. Murray
\paper Uniform analyticity of orthogonal projections
\jour  Trans. Amer. Math. Soc.
\vol   312
\yr    1989
\pages 779--817
\endref

\ref
\key   4
\by    R. R. Coifman, R. Rochberg and G. Weiss
\paper Factorization theorems for Hardy spaces in several variables
\jour  Ann. of Math.
\vol   103
\yr    1976
\pages 611--635
\endref

\ref
\key   5
\by    A. C\'ordoba
\paper The disc multiplier
\jour  Duke Math. J.
\vol   58
\yr    1989
\pages 21--29
\endref

\ref
\key      6
\by       J. Garc\'{\i}a-Cuerva and J. L. Rubio de Francia
\book     Weighted norm inequalities and related topics
\publ     North-Holland
\publaddr Amsterdam
\yr       1985
\endref

\ref
\key   7
\by    F. Marcell\'an, J. S. Dehesa and A. Ronveaux
\paper On orthogonal polynomials with perturbed recurrence relations
\jour  J. Comput. Appl. Math.
\vol   30
\yr    1990
\pages 203--212
\endref

\ref
\key   8
\by    C. J. Neugebauer
\paper Inserting $A_p$ weights
\jour  Proc. Amer. Math. Soc.
\vol   87
\yr    1983
\pages 644--648
\endref

\ref
\key   9
\by    P. Nevai and W. Van Assche
\paper Compact perturbations of orthogonal polynomials
\jour  Pacific J. Math.
\vol   153
\yr    1992
\pages 163--184
\endref

\ref
\key   10
\by    C. Segovia and J. L. Torrea
\paper Vector-valued commutators and applications
\jour  Indiana Univ. Math. J.
\vol   38
\yr    1989
\pages 959--971
\endref

\ref
\key   11
\bysame % C. Segovia and J. L. Torrea
\paper Weighted inequalities for commutators of fractional and singular
       integrals
\jour  Publicacions Matem\`atiques
\vol   35
\yr    1991
\pages 209--235
\endref

\ref
\key   12
\bysame % C. Segovia and J. L. Torrea
\paper Higher order commutators for vector-valued Calder\'on-Zygmund operators
\jour  Trans. Amer. Math. Soc.
\vol   336
\yr    1993
\pages 537--556
\endref

\ref
\key   13
\by    J. L. Varona
\paper Fourier series of functions whose Hankel transform is supported on
       $[0,1]$
\jour  Constr. Approx.
\vol   10
\yr    1994
\pages 65--75
\endref

\ref
\key      14
\by       G. N. Watson
\book     A treatise on the theory of Bessel functions
\bookinfo 2nd ed.
\publ     Cambridge Univ. Press
\publaddr Cambridge
\yr       1944
\miscnote reprinted 1995
\endref

\endRefs

\enddocument


